% META: {"id":"1SPE-EXPO-EX-002-CDP","chapitre":"1SPE-EXPONENTIELLE","type_objet":"coup_de_pouce","exercice_id":"1SPE-EXPO-EX-002","status":"approved","origin":"generated","status_history":["generated","audited","accepted","approved"],"release_acceptance":"RELEASE_OWNER_BATCH_ACCEPTANCE","acceptance_closure_digest":"sha256:9b3ccf9a81c5520fbb7e03b7b2d3e7bf2057a02834b6d908bf3be5f49f3b553f","acceptance_receipt_digest":"sha256:4fad86f0282e0fa4e0419cca1ad6379ae7af5a280c04a4a90880f90a1c044242"}
\begin{coupdepouce}{1SPE-EXPO-EX-002}
\coupDePouce{1}{Pour dériver $g(x) = 2\mathrm{e}^x$, utiliser que la dérivée d'une constante fois une fonction est la constante fois la dérivée.}
\coupDePouce{2}{L'unicité de la fonction exponentielle exige \textbf{deux} conditions : $y' = y$ \textbf{et} $y(0) = 1$. Vérifier chacune séparément.}
\end{coupdepouce}
